\uuid{cvXG}
\titre{Étude de fonction}
\auteur{bodin ; hanani ; mzari}
\organisation{exo7}
\licence{CC BY-NC-SA 4.0}
\datecreate{2020-01-20}
\niveau{L1}
\module{Analyse}
\chapitre{Fonctions circulaires et hyperboliques inverses}
\sousChapitre{Autre}
\difficulte{1}
\theme{Fonctions usuelles, étude de fonction}
\exo7qcmid{lille-1.fonctions-usuelles.009}

\contenu{
\begin{question}
\qclassification{126.00}
\qtags{motcle=étude de fonction}
Soit $f(x)=\sqrt[3]{1-x^3}$. Quelles sont les assertions vraies ?
 %$g(x)=\ln (\frac{x-1}{x+1})$.

\begin{answers}
    \good{$f$ est définie sur $\Rr$.}

    \bad{$f$ est croissante.}

    \good{$f$ est une bijection de $\Rr$ dans $\Rr$.}

    \good{L'application réciproque de $f$ est $f$.}
\end{answers}
\begin{explanations}
La fonction $x\to \sqrt[3]{x}$ est définie sue $\Rr$, elle est strictement croissante et établit une bijection de $\Rr$ dans $\Rr$. On déduit que $f$ est définie sur $\Rr$, elle est strictement décroissante et établit une bijection de $\Rr$ dans $\Rr$. Soit $y\in \Rr$, on a : $y=\sqrt[3]{1-x^3} \Leftrightarrow 1-x^3=y^3 \Leftrightarrow  x=\sqrt[3]{1-y^3}$. Donc l'application réciproque de $f$ est $f$.
\end{explanations}
\end{question}
}
